\documentclass[a4paper,11pt]{article}
\usepackage[utf8]{inputenc}
\usepackage[T1]{fontenc}
\usepackage[english]{babel}
\usepackage[left=2.5cm,right=2.5cm,top=2cm,bottom=2cm]{geometry}
\usepackage{hyperref}
\usepackage{xcolor}
\usepackage{microtype}
\usepackage{lmodern}

\definecolor{primary}{RGB}{0, 45, 114}
\hypersetup{colorlinks=true, urlcolor=primary}

\begin{document}
\noindent
\textbf{Loïc Aron MBASSI EWOLO} \\
ENSEA -- Double Diplôme Ingénieur \\
\href{mailto:loicaron.mbassiewolo@ensea.fr}{loicaron.mbassiewolo@ensea.fr} \quad (+33) 7 59 52 33 98 \\
\href{https://github.com/Nameless0l}{github.com/Nameless0l} \quad \href{https://mbassiloic.tech}{mbassiloic.tech}

\vspace{1em}
\noindent Cergy, le 22 mars 2026

\vspace{1em}
\noindent
\textbf{Objet : Candidature -- Cybersécurité IA \& Red Teaming} \\
\textbf{Référence : Offre d'alternance -- Thales La Ciotat}

\vspace{1.5em}
\noindent Madame, Monsieur,

\vspace{0.8em}
Sécuriser les systèmes critiques contre les menaces émergentes via l'IA demande une compréhension profonde des vecteurs d'attaque et une architecture défensive robuste. Vous recherchez un ingénieur capable de conjuguer NLP adversarial, cryptographie appliquée et red teaming. Mon stage en cours à INRIA Saclay (équipe TRIBE, avril-août 2026) me forme à l'analyse fine des failles de sécurité dans les applications réelles, tandis que mon expertise en systèmes sécurisés (protocoles chiffrés, RGPD by design) me positionne pour contribuer à vos défenses de prochaine génération.

\vspace{0.8em}
À INRIA Saclay, je participe à l'analyse automatisée de la cohérence et de la transparence des applications mobiles de mobilité. Ce travail inclut la détection de failles logiques via NLP des politiques de confidentialité, l'analyse statique de SDKs pour identifier les risques de fuite de données de localisation, et l'évaluation de la conformité RGPD. Cette immersion en sécurité applicative m'a enseigné comment décomposer une menace en signaux detectable : écarts entre déclarations et permissions, capacités d'inférence non-explicitées, fuites implicites. Le projet Women Safe (détection de contenus adversariaux en NLP) m'a exposé à la conception de modèles résistants au jailbreaking et aux manipulations sémantiques.

\vspace{0.8em}
Dans le projet Snappy (plateforme temps réel sécurisée), j'ai conçu une architecture de cryptage bout en bout mettant en œuvre le protocole Signal (Perfect Forward Secrecy, algorithme double ratchet) combiné au chiffrement hybride AES-256-GCM. Cette implémentation m'a exposé aux défis du chiffrement en production : gestion des clés, rotation planifiée, isolation des secrets, audit des opérations cryptographiques. Le déploiement embarqué sur le projet CoHoMa (robotique autonome militaire, Jetson Nvidia) m'a formé à l'optimisation IA en environnement contraint, où la sécurité et les perfs doivent coexister. Ces projets reflètent mon approche : sécurité et performance ne sont pas opposées, mais complémentaires.

\vspace{0.8em}
Je suis disponible dès septembre 2026 pour une alternance de 12 mois. Convaincu que la cybersécurité IA repose sur une alliance entre expertise cryptographique, analyse comportementale NLP et red teaming rigoureux, je suis prêt à relever vos défis de sécurisation des systèmes critiques.

\vspace{1.5em}
\noindent Veuillez agréer, Madame, Monsieur, l'expression de mes salutations distinguées.

\vspace{2em}
\noindent \textbf{Loïc Aron MBASSI EWOLO}

\end{document}
